\chapter{Transferencia del campo incidente sobre la esfera al plano tangente}
\label{sec:transferencia_campo_esfera_plano}

La representación de un campo definido sobre la superficie esférica en un plano tangente permite relacionar su distribución angular con una distribución espacial plana equivalente. Este proceso, conocido como \emph{transferencia de curvatura}, se compone de dos etapas: la proyección geométrica de los puntos del campo sobre la esféra al plano tangente, y la corrección en amplitud necesaria para conservar el flujo energético a lo largo de los rayos.

% ================================================================
\subsection{Proyección esférica al plano tangente}
\label{subsec:proyeccion_esfera_plano}

Sea el punto del campo sobre la esféra

\begin{equation}
\vb{Q} = R \sin\theta(\cos\phi\,\uvec{\imath} + \sin\phi\,\uvec{\jmath}) + R\cos\theta\,\uvec{k},
\label{eq:Q_vector}
\end{equation}

y el punto interior a la esféra que actúa como centro de proyección

\begin{equation}
\vb{A} = R\eta\,\uvec{k}.
\label{eq:A_vector}
\end{equation}

El punto correspondiente sobre el plano tangente \(z=-R\) se representa por

\begin{equation}
\vb{q} = r(\cos\phi\,\uvec{\imath} + \sin\phi\,\uvec{\jmath}) - R\,\uvec{k}.
\label{eq:q_vector}
\end{equation}

El rayo que une \(A\), \(Q\) y \(q\) define la condición de colinealidad

\begin{equation}
\vb{Q} - \vb{q} = \lambda \, \frac{\vb{q} - \vb{A}}{||\vb{q} - \vb{A}||},
\label{eq:colinealidad}
\end{equation}

de la cual se obtienen, resolviendo en coordenadas cilíndricas \((r,\phi)\), las relaciones

\begin{align}
r(\theta) &= R\,\frac{(1+\eta)\sin\theta}{\eta - \cos\theta}, 
\label{eq:r_de_theta}\\[6pt]
\lambda(\theta) &= -\,R\,\frac{1+\cos\theta}{|\eta-\cos\theta|}
\,\sqrt{1+\eta^2 - 2\eta\cos\theta}.
\label{eq:lambda_de_theta}
\end{align}

La primera expresión determina la correspondencia entre un punto sobre la esféra \(\theta, \phi \) y un punto sobre el plano tangente \(r(\theta), \phi\).  
La segunda describe la distancia que hay entre estos puntos $q$ y $Q$ conjugados por la proyección.  

Invirtiendo \eqref{eq:r_de_theta} se obtiene la dependencia explícita \(\theta = \theta(r)\) que se usa en el operador de transferencia:
\begin{equation}
    \theta(r)
    =
    2 \arctan\!\left(
    \frac{
        1 - \sqrt{1 - \dfrac{\eta-1}{\eta+1}\left(\dfrac{r}{R}\right)^2}
    }{\dfrac{r}{R}}
    \right),
    \label{eq:theta_de_r}
\end{equation}
donde se ha elegido la rama que satisface $\theta \to 0$ cuando $r \to 0$.

\vspace{0.5em}

% ================================================================
\subsection{Corrección en la amplitud para asegurar conservación de la energía}
\label{subsec:correccion_amplitud}

El paso de la esféra al plano debe preservar el flujo energético transportado por los rayos. Para ello se incluye un factor de oblicuidad y el jacobiano de la transformación de área.  

Sea $\uvec u$ la dirección de propagación del rayo, $\uvec{n}_S = \uvec{N}$ la normal unitaria a la esféra y $\uvec{n}_P = \uvec{k}$ la normal al plano tangente.

La conservación del flujo a lo largo de cada rayo se expresa como

\begin{equation}
|\vb{E}_P|^2 (\uvec{u}\cdot\uvec{n}_P)\, dS_P
=
|\vb{E}_S|^2 (\uvec{u}\cdot\uvec{n}_S)\, dS_S,
\label{eq:flux_conservation_sphere}
\end{equation}

donde los elementos de área están dados por

\begin{equation}
dS_P = r\,dr\,d\phi, \qquad dS_S = R^2 \sin\theta\, d\theta\, d\phi.
\end{equation}

De la ecuación \eqref{eq:flux_conservation_sphere} se deduce el factor de amplitud que asegura la conservación del flujo entre ambas superficies,

\begin{equation}
T(\theta) = \frac{|\eta - \cos\theta|}{1 + \eta}.
\label{eq:factor_amplitud}
\end{equation}

Este factor incorpora simultáneamente la proyección oblicua y el jacobiano geométrico de la transformación entre coordenadas esféricas y planas.

% ================================================================
\subsection{Operador de transferencia}
\label{subsec:operador_transferencia}

Combinando la proyección geometrica central por el rayo en dirección $\uvec u$ junto con la fase de propagación por el trayecto recorrido, dada por \(\lambda(\theta)\), y la corrección de amplitud \(T(\theta)\), se define el operador de transferencia que lleva el campo escalar \(U_S(Q) = U_S(\theta,\phi)\) definido sobre la esféra al campo \(U_P(q) = U_P(r(\theta),\phi)\) sobre el plano tangente (\(r(\theta)\) es justamente la función que caracteriza la proyección geometrica \eqref{eq:r_de_theta}):

\begin{equation}
U_P(r(\theta),\phi)
=
T(\theta)\,
e^{i k \lambda(\theta)}\,
U_S(\theta,\phi),
\qquad
r = r(\theta),
\label{eq:transferencia_campo_esfera_plano}
\end{equation}

donde \(r(\theta)\) y \(\lambda(\theta)\) están dados por las ecuaciones \eqref{eq:r_de_theta} y \eqref{eq:lambda_de_theta}, respectivamente, y \(T(\theta)\) por \eqref{eq:factor_amplitud}.  

De modo compacto, el \emph{operador de transferencia de curvatura} se escribe como

\begin{equation}
\boxed{
\mcal{M}_{S\to P}[U_S](r,\phi)
=
\frac{|\eta-\cos\theta|}{1+\eta}\,
e^{i k \lambda(\theta)}\,
U_S(\theta,\phi)
},
\label{eq:operador_transferencia_final}
\end{equation}

donde el ángulo polar \(\theta\) debe entenderse como función de \(r\) dada por \eqref{eq:theta_de_r}.  
Este operador establece la correspondencia unívoca entre la distribución angular del campo incidente sobre la esféra y su distribución espacial sobre el plano tangente, conservando el flujo y la fase óptica a lo largo de los rayos.
\begin{figure}[ht]
\centering
\begin{tikzpicture}[scale=3.1, line cap=round, line join=round, >=latex]
  %-----------------------------
  % Parámetros
  %-----------------------------
  \pgfmathsetmacro{\R}{1.25}           % radio de la esfera
  \pgfmathsetmacro{\eta}{1.0/1.6}      % A = R*eta
  \pgfmathsetmacro{\thp}{40}           % theta' (grados) en [0,90]
  % theta = 180 - theta'
  \pgfmathsetmacro{\costh}{-cos(\thp)} % cos(theta)
  \pgfmathsetmacro{\sinth}{ sin(\thp)} % sin(theta)
  %-----------------------------
  % Coordenadas
  %-----------------------------
  \pgfmathsetmacro{\Qx}{\R*\costh}
  \pgfmathsetmacro{\Qy}{\R*\sinth}
  \pgfmathsetmacro{\qy}{\R*(1+\eta)*\sinth/(\eta-\costh)}
  \pgfmathsetmacro{\Ax}{\R*\eta}
  \pgfmathsetmacro{\Ay}{0}
  \pgfmathsetmacro{\kx}{\Qx-\Ax}
  \pgfmathsetmacro{\ky}{\Qy-\Ay}
  \pgfmathsetmacro{\k}{0.6}
  \pgfmathsetmacro{\Px}{\Qx+\k*\kx}
  \pgfmathsetmacro{\Py}{\Qy+\k*\ky}
  \coordinate (C) at (0,0);
  \coordinate (A) at (\Ax,\Ay);
  \coordinate (Q) at (\Qx,\Qy);
  \coordinate (q) at ({-\R},\qy);
  \coordinate (Pout) at (\Px,\Py);
  %-----------------------------
  % Plano z = -R
  %-----------------------------
  \draw[line width=1pt, black!80] (-\R,-1.35*\R) -- (-\R,1.35*\R);
  \node[rotate=90, anchor=south] at (-\R,-0.45*\R) {Plano $z{=}{-}R$};
  %-----------------------------
  % Eje z
  %-----------------------------
  \draw[->, line width=0.7pt, black!70] (-1.4*\R,0) -- (1.45*\R,0) node[right] {$z$};
  %-----------------------------
  % Esfera
  %-----------------------------
  \draw[line width=0.8pt, black!80] (C) circle (\R);
  %-----------------------------
  % Radio C->Q (tenue)
  %-----------------------------
  \draw[line width=0.4pt, black!80, densely dotted] (C)--(Q);
  %-----------------------------
  % Ángulo theta'
  %-----------------------------
  \pgfmathsetmacro{\rtheta}{0.4}
  \draw[line width=0.6pt, black!70] ({-\rtheta},0) arc (180:{180-\thp}:\rtheta);
  \node at ({1.6*\rtheta*cos(180 - \thp/2)},{1.6*\rtheta*sin(180 - \thp/2)}) {$\theta'$};
  %-----------------------------
  % Puntos
  %-----------------------------
  \fill (C) circle (0.4pt) node[below=3pt]{$C$};
  \fill (A) circle (0.5pt) node[below right=1pt]{$A$};
  \fill (Q) circle (0.5pt) node[below=5pt]{$Q$};
  \fill (q) circle (0.5pt) node[below=4pt, left=1pt]{$q$};
  %-----------------------------
  % Rayo exterior
  %-----------------------------
  \draw[line width=1.2pt, ->, black] (Pout)--(Q) node[pos=0.25, above left, sloped] {$\uvec{u}$};
  %-----------------------------
  % Prolongación interior
  %-----------------------------
  \draw[line width=0.8pt, densely dashed] (Q)--(A);
  %-----------------------------
  % Normal esférica n_S en Q (color azul)
  %-----------------------------
  \pgfmathsetmacro{\nslength}{0.18}
  \draw[->, line width=0.7pt, blue!70!black] (Q)--++({(\nslength)*\costh},{(\nslength)*\sinth}) 
        node[pos=0, above=14pt, right=-9pt]{$\hat{n}_S$};
  %-----------------------------
  % Normal al plano n_P en q (color rojo)
  %-----------------------------
  \pgfmathsetmacro{\nplength}{0.18}
  \draw[->, line width=0.7pt, red!70!black] (q)--++(\nplength,0) 
        node[pos=0, above=8pt, right=-3pt]{$\hat{n}_P$};
\end{tikzpicture}
\caption{Proyección al plano tangente con \(z\) horizontal y uso de \(\theta' \in [0,\pi/2]\) (\(\theta=\pi-\theta'\)).
El rayo llega desde el exterior (línea continua) hasta \(Q\) y se prolonga (línea discontinua) hacia $A=R(n_0/n_i)$. El plano \(z=-R\) es la recta vertical a la izquierda.}
\label{fig:sphere_to_plane_projection_final}
\end{figure}
